\chapter*{Glosario}
\addcontentsline{toc}{chapter}{Glosario}

\begin{description}
\item[Broker] Intermediario de MQTT. Recibe los mensajes de quienes publican y los entrega a quienes se
suscribieron. Ningún sensor habla directamente con la plataforma: todos hablan con el broker.
\item[Checksum (suma de comprobación)] Dígito de verificación al final del mensaje Contact ID. Si un
dígito se corrompe en el camino, la suma no cuadra y el mensaje se rechaza.
\item[ChirpStack] Servidor de red LoRaWAN de código abierto. Es el «cerebro» de la red de radio:
descifra los mensajes de los nodos, elimina duplicados, administra los dispositivos y entrega el dato
limpio a la plataforma mediante MQTT. El gateway, en cambio, solo reenvía sin entender.
\item[Contact ID (SIA DC-05)] Protocolo de 16 dígitos de la industria de alarmas para reportar eventos.
DC-05 es el nombre del estándar; Contact ID, el del protocolo (como IEEE 802.11 y WiFi).
\item[Datos extendidos] Campo \texttt{[<A--Z><datos>]} que el estándar DC-09 reserva para información
definida por el implementador. En este trabajo transporta la medición junto al evento.
\item[DC-09] Estándar SIA que transporta mensajes Contact ID por internet (TCP/UDP) con cifrado AES.
\item[DTMF] Tonos del teclado telefónico. Fue el transporte original de Contact ID.
\item[Gateway] Antena que recibe la radio de los nodos LoRa y la reenvía por internet al servidor de
red. No posee las claves de cifrado: es un traductor sin capacidad de leer el contenido.
\item[LoRa] \emph{Long Range}. Técnica de modulación de radio de largo alcance y bajo consumo (capa
física). Son ondas de radio como las de la telefonía celular; la diferencia está en la banda (915\,MHz
de uso libre, sin licencia ni operador) y en el protocolo que viaja sobre ellas.
\item[LoRaWAN] \emph{Long Range Wide Area Network}. Protocolo de red construido sobre LoRa: registro de dispositivos, cifrado, entrega.
\item[MQTT] \emph{Message Queuing Telemetry Transport}. Protocolo de mensajería «publicar/suscribir»: los sensores publican en un tópico
(\texttt{sitio/invernadero-2/sensor-007/temperatura}) y cualquier programa suscrito lo recibe, sin que
emisor y receptor se conozcan. El dispositivo inicia la conexión saliente, por lo que funciona detrás
de cualquier router sin abrir puertos.
\item[NAT y puertos] El router de una red doméstica o empresarial comparte una sola dirección pública
entre todos sus dispositivos (\emph{Network Address Translation}). Deja salir cualquier conexión
iniciada desde dentro y recuerda a quién devolver la respuesta, pero \textbf{bloquea toda conexión que
llegue desde fuera} sin que nadie la haya pedido. Para que un servidor externo pudiera «llamar» a un
sensor, habría que abrir un puerto en el router y redirigirlo al sensor (\emph{port forwarding}): una
configuración manual por cada red, y una puerta expuesta a internet. Si en cambio el sensor inicia la
conexión hacia la plataforma, el router la deja pasar sin configurar nada. Es la razón del
requerimiento RNF-10.
\item[Payload] Los bytes útiles de un mensaje, descontadas las cabeceras del transporte. En LoRaWAN va
de 11 bytes (tasa más lenta, máximo alcance) a 51--53 bytes.
\item[PPDIOO] Metodología de ciclo de vida de redes de Cisco: Prepare, Plan, Design, Implement,
Operate, Optimize.
\item[RSSI] \emph{Received Signal Strength Indicator}: cuánta potencia de señal llegó al receptor, en
dBm (típicamente entre $-30$, excelente, y $-120$, casi inaudible). Indica \emph{qué tan fuerte} se oye
al nodo.
\item[SNR] \emph{Signal-to-Noise Ratio}: cuánto sobresale la señal por encima del ruido de fondo, en dB.
Indica \emph{qué tan limpia} se oye. Un nodo puede tener RSSI aceptable y SNR malo si hay
interferencia: ambos juntos describen la salud del enlace.
\item[TimescaleDB] Extensión de PostgreSQL para datos con marca temporal. Sigue siendo una base de datos
relacional con SQL estándar.
\item[Tópico] La «ruta» de un mensaje MQTT, con forma de carpetas separadas por barras.
\end{description}
